\chapter{Einleitung}
\label{chap:einleitung}

\section{Motivation und Problemstellung}
\label{sec:motivation}

Der digitale Handel hat sich in den vergangenen Jahren zu einem
zentralen Bestandteil des wirtschaftlichen Lebens entwickelt. Laut dem
Bundesverband E-Commerce und Versandhandel Deutschland erzielte der
deutsche E-Commerce-Markt im Jahr 2024 einen Bruttoumsatz mit Waren
von rund 80,6 Milliarden Euro \cite{bevh_2025}. Deutschland zählt zu
den größten E-Commerce-Märkten Europas
\cite{trade_gov_germany_ecommerce}.

Der Umsatz verteilt sich dabei ungleich. Auf Marktplätze allein
entfielen 44 Milliarden Euro, während Pure Player,
Multichannel-Händler und D2C-Anbieter zurückgingen \cite{bevh_2025};
unter den einzelnen Anbietern führen amazon.de und otto.de die
Rangliste an, mit Nettoumsätzen von 15,76 beziehungsweise 4,54
Milliarden US-Dollar \cite{trade_gov_germany_ecommerce}. Der Markt
konzentriert sich damit auf wenige große Plattformen, die über
erhebliche technische und finanzielle Ressourcen verfügen.

Insbesondere kleine und mittelständische Unternehmen stehen dabei vor
einer Hürde: Die Einführung digitaler Lösungen, ausdrücklich auch von
E-Commerce-Systemen, verlangt Investitionen in Hardware, Software und
Schulung, die für kleinere Unternehmen unerschwinglich sein können. In
einer Erhebung unter 4531 kolumbianischen KMU zählen hohe
Investitionskosten branchenübergreifend zu den drei am stärksten
bewerteten Hürden der Digitalisierung
\cite{digitalization_barrier_smes}. Hinzu kommt, dass bestehende
Plattformen häufig entweder den Produkthandel oder die Buchung von
Dienstleistungen abdecken, selten jedoch beides in einem einheitlichen
System. So beschränken sich Amazon und Etsy im Wesentlichen auf den
Produkthandel, während Booking.com und Fiverr ausschließlich
Dienstleistungen vermitteln; eine systematische Gegenüberstellung
dieser Plattformen anhand von vier Merkmalen findet sich in
Abschnitt~\ref{subsec:wettbewerbsanalyse}.

Aus dieser Marktlücke entstand die Idee zur Plattform
\textit{Bazario}, die bei SahemTech entwickelt wird. \textit{Bazario}
verfolgt das Ziel, Händlern und Dienstleistern eine gemeinsame
digitale Infrastruktur bereitzustellen, auf der Produkte verkauft und
Dienstleistungen gebucht werden können, ohne dass die Anbieter eigene
Systeme betreiben müssen. Gleichzeitig ermöglicht die Plattform
Käufern, beide Angebotsformen in einem einzigen Bestellprozess zu
kombinieren.

Die technische Umsetzung einer solchen Plattform stellt hohe
Anforderungen an das Frontend. Es muss drei Nutzerrollen mit jeweils
eigenen Funktionsbereichen und Zugriffsrechten bedienen, nämlich
Käufer, Verkäufer und Dienstleister; die vierte Rolle, der
Administrator, greift über das vom Backend bereitgestellte Dashboard
auf die Plattform zu. Im weiteren Verlauf der Arbeit werden für diese
Rollen die plattforminternen Bezeichnungen \textit{Customer},
\textit{Seller}, \textit{ServiceProvider} und \textit{Admin}
verwendet. Gleichzeitig soll die Oberfläche konsistent, intuitiv und
performant sein; was darunter im Einzelnen zu verstehen ist, legt
Abschnitt~\ref{subsec:nichtfunktionale-anforderungen} in messbaren
Größen fest. Die Entwicklung eines solchen Multi-Rollen-Frontends
bildet den Kern dieser Arbeit.

\section{Ziel der Arbeit}
\label{sec:ziel}

Ziel dieser Bachelorarbeit ist der Entwurf und die prototypische
Entwicklung des Frontends der Webanwendung \textit{Bazario}, einer
Multi-Rollen-Plattform, die E-Commerce und Dienstleistungsbuchungen in
einem einheitlichen System zusammenführt. Im Fokus steht die
Konzeption einer rollenbasierten Single-Page Application, die alle
relevanten Funktionsbereiche der Plattform abdeckt und eine
konsistente, benutzerfreundliche Oberfläche für die drei im Frontend
abgebildeten Nutzerrollen bereitstellt.

Das Frontend bezieht seine Daten über die vom Backend bereitgestellten
\acfirst{REST}-API-Schnittstellen; allein der Nachrichtenbereich nutzt
zusätzlich eine dauerhafte Verbindung über einen Echtzeitdienst. Zum
Einsatz kommen dabei moderne Webtechnologien wie React, die Bibliothek
\textit{Zustand}, React Router und Tailwind \acfirst{CSS}.

Die gewonnenen Erkenntnisse münden in einen lauffähigen Prototypen,
der die zentralen Nutzerflows aller drei Rollen abbildet und anhand
funktionaler sowie nicht-funktionaler Kriterien bewertet wird.

\section{Abgrenzung}
\label{sec:abgrenzung}

Die vorliegende Arbeit beschränkt sich ausschließlich auf die
Entwicklung des Frontends von \textit{Bazario}. Das zugehörige Backend wurde im Rahmen einer begleitenden
Bachelorarbeit konzipiert und umgesetzt \cite{alaatki_backend}. Es basiert auf dem
Laravel-Framework, verwendet eine relationale MySQL-Datenbank,
setzt Laravel Sanctum zur tokenbasierten Authentifizierung ein
und nutzt Pusher als WebSocket-Dienst für die
Echtzeitkommunikation. Die Zahlungsabwicklung erfolgt
serverseitig über Stripe einschließlich Webhook-Verarbeitung.

Serverseitige Logik, Datenbankstruktur, Deployment und Hosting
sind damit explizit nicht Gegenstand dieser Arbeit. Die
Frontend-Implementierung baut auf den definierten
\acs{REST}-API-Schnittstellen des Backends auf und setzt deren
korrekte Funktionsweise voraus. Ebenso sind Aspekte wie
Suchmaschinenoptimierung, serverseitiges Rendering sowie
die Entwicklung einer mobilen nativen Anwendung nicht Teil
dieser Arbeit.

Die Plattform ist mehrsprachig ausgelegt. Dabei sind zwei Arten von
Texten zu unterscheiden: die Beschriftungen und Meldungen der
Oberfläche, die im Frontend liegen, und die inhaltlichen Angaben zu
Produkten, Dienstleistungen und Anbietern, die das Backend liefert.
Die Oberfläche ist Gegenstand dieser Arbeit und in Deutsch, Englisch
und Arabisch umgesetzt; für Arabisch wechselt zusätzlich die
Schreibrichtung des Dokuments. Die inhaltlichen Angaben werden vom
Backend in der angeforderten Sprache ausgeliefert, wobei die Formulare
der Anbieter Titel und Beschreibung in Englisch und Arabisch erfassen. Welche
Sprachfassungen eines Inhalts entstehen, ist damit eine Frage des
Datenmodells und nicht der Sprachwahl in der Oberfläche.

Die administrativen Funktionen der Plattform, etwa die Prüfung
von Rollenwechselanträgen und die Freigabe von Werbeanzeigen,
erfolgen über das vom Backend bereitgestellte \textit{Admin}-Dashboard
und sind nicht Teil des in dieser Arbeit entwickelten Frontends.

\section{Aufbau der Arbeit}
\label{sec:aufbau}

Die Arbeit gliedert sich in sieben Kapitel. Nach dieser
Einleitung vermittelt Kapitel~\ref{chap:theoretischer-hintergrund}
die theoretischen Grundlagen zu E-Commerce, Marktplatzmodellen und
Frontend-Architektur und stellt die für eine solche Anwendung in
Betracht kommenden Frameworks und Kommunikationsansätze
gegenüber. Kapitel~\ref{chap:ist_analyse} analysiert den
Ausgangszustand des Projekts, leitet die funktionalen sowie
nicht-funktionalen Anforderungen an das Frontend ab und begründet
auf dieser Grundlage die Auswahl des eingesetzten Technologiestacks.

Kapitel~\ref{chap:konzept} stellt das erarbeitete
Frontend-Konzept vor, einschließlich Architekturprinzipien,
\acfirst{UI}/\acfirst{UX}-Konzept und der Planung der einzelnen Funktionsbereiche.
Die konkrete Umsetzung ausgewählter Komponenten und Module
wird in Kapitel~\ref{chap:umsetzung} anhand repräsentativer
Implementierungsbeispiele dokumentiert. Kapitel~\ref{chap:testen}
beschreibt die angewandte Teststrategie und bewertet die
Ergebnisse anhand funktionaler und nicht-funktionaler Kriterien.
Abschließend fasst Kapitel~\ref{chap:fazit} die wichtigsten
Erkenntnisse zusammen und gibt einen Ausblick auf mögliche
Erweiterungen und technische Weiterentwicklungen der Plattform.